\subsection{8.3  Multi-Agent Mesh Deployments: The New Hire Problem}

The APIP framework as presented assumes a single-agent deployment with a clear locus of accountability: one agent, one behavioral contract, one remediation process. This assumption breaks down in multi-agent systems---increasingly common in enterprise deployments---where a cluster of agents cooperates to handle a composite task, sharing context, tools, memory, and sometimes overlapping task surfaces.

In such agent meshes, a class of failure mode emerges that is structurally invisible to the single-agent APIP: mesh disruption. Consider the introduction of a new specialist agent---Agent N---into an existing mesh. Agent N performs its designated task well. Yet aggregate system performance, or the individual performance of peer agents, declines. The APIP framework triggers on the peer agents, whose metrics have crossed threshold. Cause attribution, however, fails to find an intrinsic failure mode: the peer agents have not drifted, their retrieval is fresh, their alignment is intact. The root cause is exogenous---and the single-agent APIP has no mechanism to represent or address it.

Kim et al. [28] provide direct empirical support: their large controlled study of multi-agent scaling found performance change relative to single-agent baselines ranging from +80.8\% on decomposable tasks to \textbackslash{}u221270.0\% on sequential planning, identifying a coordination-saturation effect where additional agents suppress rather than improve reasoning. Introducing Agent N into an existing mesh is precisely this kind of architecture change\textbackslash{}u2014and the degradation manifests in peer agents, not in Agent N itself, making standard per-agent monitoring blind to the cause. This maps onto a well-documented class of phenomena in organizational behavior. Morgeson, Mitchell, and Liu's [16] Event System Theory provides the overarching framework: the introduction of Agent N is an organizational event whose impact is a function of its novelty (the mesh has not previously included this capability), disruptiveness (it changes the shared context distribution), and criticality (it affects the mesh's output quality). Strizver and Ployhart [19] extended EST specifically to team disruption and argued that all changes in team composition, emergent states, and structure only matter for team disruption to the extent they influence coordination---a claim that maps directly onto the agent mesh case, where the disruption manifests as coordination failure in shared context.

We hypothesize four mechanisms by which a locally high-performing agent can produce performance externalities in its mesh peers, each grounded in organizational research:

Context Space Crowding. If agents share a context window or working memory buffer, Agent N's high-confidence, well-formed outputs may dominate the shared context, displacing signals that peer agents depend on---an attention competition problem at the systems level. This maps to the Transactive Memory Systems literature [23, 15, 18]: when Agent N enters and changes the content of shared memory, the existing agents' TMS---their calibration to what's in the shared context and who produced it---is disrupted. Ren and Argote [18] document that TMS disruption from membership change degrades coordination. From a systems perspective, this is the coherence failure described in multi-agent memory research [29]: without a coordination protocol governing shared context writes, one agent\textbackslash{}u2019s outputs overwrite or displace another\textbackslash{}u2019s.

Role Boundary Erosion. Agent N, performing its task thoroughly, may produce outputs that functionally overlap with the input surface of adjacent agents, causing those agents to receive redundant, contradictory, or already-processed inputs. Summers, Humphrey, and Ferris [20] found that membership changes involving strategically core roles produced the highest levels of coordination flux---precisely the dynamic at play when Agent N occupies a central position in the mesh's task graph.

Distributional Shift in Shared State. Agents calibrated against a certain content distribution in shared memory now operate out-of-distribution when Agent N begins populating that space. Lewis, Belliveau, Herndon, and Keller [15] demonstrated that group cognition is disrupted by membership change in part because the collective knowledge structure shifts---existing members' expectations about the group's knowledge no longer hold.

Trust Calibration Collapse. In architectures where agents implicitly weight peer outputs, a high-performing Agent N may receive disproportionate deference from the orchestrator or peer agents, causing valid contributions from other agents to be systematically underweighted. This mirrors the organizational finding that newcomer attributes interact with team norms to reshape coordination patterns [22].

This problem is directly analogous to a well-documented organizational phenomenon: the new hire who is individually excellent but systemically disruptive---who reshapes team communication patterns, crowds out peer contributions, or shifts implicit norms in ways that reduce collective performance. The newcomer socialization literature documents that support for newcomers from coworkers declines within the first 90 days [6], and that structured onboarding---including socialization tactics, buddy systems, and norm communication [2, 22]---significantly mitigates the disruptiveness of the transition process.

From the APIP perspective, mesh disruption creates a causal attribution failure: the underperforming agents' APIPs cannot be resolved because the root cause is not remediable at the agent level. This motivates two extensions to the framework. First, a mesh-aware cause attribution category---exogenous disruption---should be added to the failure mode taxonomy (reflected in the schema in Table 2), with a corresponding diagnosis protocol that includes analysis of shared context composition and peer agent behavioral change events. Second, a mesh onboarding protocol is needed: a structured introduction process for new agents that includes context scoping (isolating the new agent's context contributions during a probationary window), behavioral contract alignment across the mesh, and regression monitoring on peer agents during the introduction period---analogous to the 30-60-90 day structured onboarding plans documented in the newcomer socialization literature [2, 6].

We identify the formal treatment of agent mesh disruption as a significant open problem at the intersection of multi-agent coordination, information-theoretic resource allocation, and AI governance. The mechanism design framing is particularly productive: mesh disruption can be modeled as a congestion externality in shared context space, with the APIP framework providing the remediation layer once the externality is detected. Krishnan [30] notes that even with standardized context protocols such as MCP, tracing causality across hundreds of agents exceeds current tooling---confirming that the attribution failure we describe is a documented practical constraint. A full formal treatment is reserved for a companion paper.

